% ===================================================================
% limitations.tex -- Limitations of the computational study (data, device model, first-principles calculations,
% procedure). Draft 2026-10-13 (Claude), for human review; no claim here is approved for a manuscript (charter section 9).
% ASCII only. Every number carries a source comment  % <!-- KEY:line -->  (file:line). KEY legend:
%   PRIV    = iwo-tft-research-private/                     PKG    = IWO_PHYSICS_CONSTRAINED_MODEL_V1/
%   CHARTER = PRIV/00_PROJECT_CHARTER.md                    CS     = PRIV/01_CURRENT_STATE.md
%   LEDGER  = PRIV/03_CLAIMS_LEDGER.yaml                    DLOG   = PRIV/05_DECISION_LOG.md
%   REG     = PRIV/codex/work_chat/HYPOTHESIS_REGISTER.md   SLAB   = PRIV/compute/amorph/SLAB_PROTOCOL.md
%   MLV     = PRIV/claude/plans/ML_VALIDATION.md            V3     = PKG/config/iwo_material_model_v3.yaml
%   CD      = PKG/analysis_2026-10-10_campaign/CAMPAIGN_DESIGN.md
%   S2PREP  = PKG/analysis_2026-10-10_campaign/S2_PREP.md   S2FIX  = PKG/analysis_2026-10-10_campaign/S2_SCORER_FIXES_2026-10-12.md
%   V2TEX   = PKG/report_latex_full/chapters/14_model_v2.tex
% ===================================================================
\label{sec:limitations}

\paragraph{Measured data.}
\begin{itemize}
\item One device per thickness, one drain bias (\SI{0.7}{V}), no gate current and no reverse sweep: device-to-device
spread cannot be quoted, and no thickness trend can be separated from it. % <!-- LEDGER:5-9 ; LEDGER:26 -->
\item The thicknesses are nominal, and whether the devices come from one fabrication run, and their sweep history,
are unknown. % <!-- CS:253-254 -->
\item The \SI{2.0}{nm} off-state floor (\SI{1.3}{pA} in total) is at the level of the measurement set-up. % <!-- LEDGER:20 -->
\item The \SI{31.8}{nm} device never turns off; it is a documented model failure and not a holdout. % <!-- LEDGER:172-176 ; CS:245-246 -->
\item The source paper's mobility and near-thermal swing differ from ours on the same curves; a definition difference
reconciles the mobility numerically but the paper's recipe is unconfirmed (claims C024, C025 \prov{UNRESOLVED}). % <!-- LEDGER:138-152 -->
\item No new measurements were possible (decision D-11). The measurement that separates an ohmic parallel strip
(floor $\times7$ from 0.1 to \SI{0.7}{V}) from a diffusion form ($\times1.02$) could not be made, so the TCAD floor
tests bound and rank mechanisms only within the model. % <!-- DLOG:22 ; CS:219-220 ; CD:395 -->
\item Which oxide touches the channel is not confirmed by the fabrication group; the model places \AlO{} against the
channel. % <!-- CS:241 ; SLAB:68-70 -->
\item Every comparison with device data is retrospective; leave-one-film-out tests measure transferability within four
already-seen devices, not independent prediction. % <!-- REG:37 ; CD:501 -->
\end{itemize}

\paragraph{Device model.}
\begin{itemize}
\item The parameters are not unique: work function, affinity, fixed charge and conduction-band shift act as one number,
the band mobility trades off against the tail density, and series resistance against the mobility exponent (claim
C018). Confinement is physically expected but cannot be identified from transfer curves. % <!-- LEDGER:104-107 ; LEDGER:118 ; CS:153-155 -->
\item Every out-of-sample test of the \SI{6.3}{nm} threshold has failed (claim C014, \prov{FALSIFIED}); the earlier
models are calibrated descriptions, not predictive models of the thickness dependence. % <!-- LEDGER:80-84 ; V2TEX:334-335 -->
\item The earlier calibrated model did not reproduce the gate-voltage dependence of $\muFE$ at 2.0 and \SI{6.3}{nm}, and
its \SI{13.2}{nm} swing is too large (claims C012, C010). % <!-- LEDGER:57-61 ; LEDGER:69-72 -->
\item The quasi-two-dimensional density of states at \SI{2.0}{nm} is a substitution inside the 3-D Fermi--Dirac form; it
overstates the degenerate on-state density by up to $1.7\times$, and the on-current at \SI{3}{V} by an inferred
7--8\,\%. The density-of-states convention cannot be told apart from fixed charge, tail density and Gaussian depth in
subthreshold data. % <!-- CD:111 ; REG:226-228 -->
\item The rigid-shift interpretation of the confinement tests holds only because trap levels follow $\Ec$ in the model
(\prov{ASSUMED}). % <!-- REG:206 -->
\item The DFT laws are extrapolated beyond \SI{2}{nm}, the HSE06/PBE ratio was measured at \SI{1}{nm} only, and the
device thickness is identified with the slab H--H thickness (\prov{ASSUMED}). % <!-- CD:151 ; V3:76 -->
\item Contacts are ideal Ohmic and unconstrained (no transfer-length data); the solution is electron-only; the donor
density is constant by assumption. % <!-- V3:209-211 ; V3:124-125 -->
\item The leakage-maximising oxide set is an optimistic scenario, not a bound: pinholes, local thinning and percolation
paths are not modelled. The oxide trap-assisted tunnelling runs are not convergence-qualified. % <!-- CD:376 ; REG:157 ; S2PREP:164 -->
\item The numerical-resolution method of the scorer is an \prov{ASSUMED} operational method; with the current run plan
most primary leakage and affinity-only verdicts will be provisional or unresolved. % <!-- S2FIX:238-247 ; S2FIX:288-291 -->
\end{itemize}

\paragraph{First-principles calculations.}
\begin{itemize}
\item PBE without a Hubbard correction, spin-unpolarised and neutral cells: the results can show consistency with
donor behaviour, not full ionisation or charge-transition levels. Hybrid functionals enter only as checks, and HSE06
itself underestimates the \InO{} gap. % <!-- REG:260-262 ; REG:585 ; REG:544 -->
\item The absolute gap is anchored to a measurement on a different sputtered IWO device. % <!-- REG:541 -->
\item The machine-learned structures failed the pre-registered force check on the W first shell and are DFT starting
points, not validated amorphous models (claim C037). % <!-- LEDGER:228-233 -->
\item One melt--quench sample per structure (two at \SI{2.0}{nm}); two seeds give diagnostic variability, not a
confidence interval, and a difference smaller than the gap between the two \SI{2.0}{nm} samples is not a result.
Force agreement on a few snapshots would not show that the ensemble represents sputtered films. % <!-- MLV:35-36 ; REG:282 -->
\item Small cells: W images at least \SI{12.5}{\angstrom} apart are only partly controlled, and one excess electron per W
gives about $6\times10^{20}$\cmcu{}, far above the channel carrier density. % <!-- REG:622 ; SLAB:126 -->
\item One state per DFT cell corresponds to about $6\times10^{13}$\cmsq{}, so the fitted fixed charge and interface
trap density of the device model cannot be tested by DFT. % <!-- REG:40 -->
\item Mixed relaxation-smearing families in the W series; a provisional final-SCF width without an electron-affinity
comparison; restart non-reproducibility at production tolerance in one cell (claim C036). % <!-- REG:939-943 ; REG:1025 ; REG:670 ; LEDGER:221-226 -->
\item Deadline effects: relaxations not converged at the cutoff are exploratory; the \SI{1.5}{nm} cells were cut, the
\SI{2.0}{nm} \AlO{} interface cell was stopped and the \SI{2.0}{nm} bare control parked, so H-I is expected to be
unresolved and H-H unresolved at \SI{2.0}{nm}. % <!-- REG:1076 ; DLOG:46 ; REG:823-827 -->
\item H-N is exploratory, and its \SI{2.0}{nm} endpoint is unresolved without a certified C4 polish of both seeds. % <!-- DLOG:39 ; REG:1016-1018 -->
\item The amorphous effective-mass workflow was descoped, so the mass-only mobility prediction is not computed. % <!-- REG:170-171 -->
\item The \HfO{} cell is one seed with a different oxide and bonding; it cannot establish the stack order, and its
oxide density and possible partial crystallisation are assumptions to be checked by the registered flags. % <!-- REG:391-397 ; SLAB:141 ; SLAB:185 -->
\item The sub-nanometre crystalline slab results depend on the termination recipe and are preliminary (claim C021). % <!-- LEDGER:120-124 -->
\item Fixed-geometry hybrid results support electronic robustness only; equilibrium self-trapping would need a
structural qualification, which is not claimed. % <!-- REG:452 -->
\end{itemize}

\paragraph{Procedure.}
\begin{itemize}
\item Partial electronic unblinding of four \SI{2.0}{nm} cells (A36 erratum) and a procedural post-hoc disclosure (A35)
are part of the record (\cref{sec:m-reg-sight}). % <!-- REG:1027 ; REG:950-959 -->
\item TCAD protocol deviations: the S0-02 repair cap and launches while the regression gate was red. % <!-- REG:39 ; S2PREP:82 -->
\item AI reviewers share training priors and can share blind spots; holdout custody, external timestamps and human
judgement are the real safeguards. In this work there is no holdout and the external timestamp is pending. % <!-- CHARTER:84-85 -->
\pdPending{P31}{TEXT PENDING: final wording once the external timestamp (P27) and the charter adoption (P26) are settled}.
\end{itemize}
